\section*{Введение}
\addcontentsline{toc}{section}{Введение}

\textbf{Актуальность темы.} В условиях стремительного роста рынка электронной 
коммерции эффективность работы служб экспресс-доставки напрямую зависит от 
качества логистического планирования. Традиционные методы ручного распределения 
заказов приводят к неоптимальному использованию транспортных ресурсов, увеличению 
времени ожидания для клиентов и росту операционных издержек. Разработка 
специализированных веб-приложений, интегрирующих современные алгоритмы оптимизации 
маршрутов и средства геоинформационного мониторинга, является необходимым условием 
для повышения конкурентоспособности логистических компаний.

\textbf{Степень разработанности проблемы.} Задача маршрутизации транспортных средств 
(VRP) классифицируется в дискретной математике как NP-трудная. Основы её решения 
заложены в работах Данцига, Рамзера и Кларка. В области веб-ориентированных систем 
существенный вклад внесли современные исследователи ГИС-технологий. Несмотря на 
наличие коммерческих продуктов (напр. Яндекс.Маршрутизация), создание гибких 
open-source решений для малого и среднего бизнеса остается актуальной и не до 
конца решенной практической задачей.

\textbf{Объект исследования} — процессы организации и управления логистическими 
цепочками в службах экспресс-доставки.

\textbf{Предмет исследования} — программные средства и алгоритмические решения 
для автоматизации построения оптимальных маршрутов в рамках веб-ориентированной 
платформы.

\textbf{Целью данной работы} является проектирование и реализация клиент-серверного 
веб-приложения, предназначенного для автоматизации учета заказов и минимизации 
транспортных издержек за счет оптимизации логистических путей.
